\documentclass{article}

% Language setting
\usepackage[english]{babel}

% Page size and margins
\usepackage[letterpaper,top=2cm,bottom=2cm,left=3cm,right=3cm,marginparwidth=1.75cm]{geometry}

% Useful packages
\usepackage{amsmath,amssymb}
\usepackage{float}
\usepackage[ruled,vlined]{algorithm2e}
\usepackage{graphicx}
\usepackage[colorlinks=true, allcolors=blue]{hyperref}
\usepackage{url}
\title{Information Gain for AI-Mediated Discovery: A Content Recommendation and Evaluation Architecture for GEO Optimization}
\date{}

\begin{document}
\maketitle

\section{Introduction and Objectives}

Large Language Models (LLMs), AI-assisted search systems, and agentic answer interfaces increasingly serve as the first interface through which users explore products, compare tools, and learn technical concepts. In these settings, answers are synthesized from vendor documentation, landing pages, blogs, community pages, media, and competitor content. For organizations that invest in owned content, two practical questions follow: how content should be written so that AI systems surface it and describe it accurately, and how such behavior should be monitored and improved over time.

This project studies \emph{Information Gain} (IG) as an operational measure of content quality for AI-mediated discovery and proposes an implementable recommendation workflow for GEO optimization. A page has high IG when it contributes concrete, specific, structured, and checkable information on a topic that is not already common in comparable pages. The central hypothesis is that improving IG in brand-owned content increases the likelihood that AI systems cite the brand, mention it earlier, and reproduce the intended product narrative more faithfully. The operational goal is to score pages, identify content gaps, suggest targeted revisions, support controlled variant deployment, and re-measure AI-side outcomes.

The project evaluates three classes of outcomes:
\begin{itemize}
    \item \textbf{visibility outcomes}, such as whether the brand appears in an answer and how prominently it appears;
    \item \textbf{narrative outcomes}, such as correctness and completeness of the description of the brand offering;
    \item \textbf{site outcomes}, such as agentic referrals, visits, and selected downstream actions on optimized pages, measured through a join between variant-aware page delivery logs, referral or session instrumentation, and downstream analytics events when available.
\end{itemize}

The working causal chain is:
\[
\text{Higher IG in brand content}
\;\Rightarrow\;
\text{better AI visibility and narrative alignment}
\;\Rightarrow\;
\text{improved downstream page outcomes}.
\]

The main technical artifacts are:
\begin{itemize}
    \item a two-level IG scoring function that scores editable page snapshots and aggregates them into topic-cluster outcomes;
    \item an \emph{Information Gain Card} (IGC) for each topic cluster;
    \item a \emph{Query--Answer Trace} (QAT) representation for answer collection and metric extraction;
    \item a failure taxonomy for cases where improved content does not yield the expected runtime behavior;
    \item a hybrid human + LLM judging pipeline for scalable annotation and scoring;
    \item a recommendation output layer that translates low-scoring dimensions and runtime failures into concrete content-side actions.
\end{itemize}

\section{Conceptual Model and Abstractions}

\subsection{Information Gain as Page-Level Scoring and Cluster-Level Aggregation}
\label{sec:igscores}

A page snapshot is represented as an evidence object
\[
p = \langle u, x, m, c \rangle,
\]
where \(u\) is the canonical URL, \(x\) is the extracted textual content, \(m\) is page metadata, and \(c\) is the assigned topic cluster. The page is the unit of capture, indexing approximation, and content intervention. A topic cluster is the aggregation unit used to represent a broader user intent, combine related pages, evaluate cross-page consistency and coverage, and attach runtime QAT outcomes.

Topic clusters \(C\) correspond to recurring user intents and comparable content surfaces within the brand site. For the \textbf{Frescopa} deployment, representative examples include:
\begin{itemize}
    \item \(C_1\): \textit{coffee subscriptions and recurring delivery}, covering recurring delivery, plan selection, delivery cadence, pause or cancellation conditions, and suitability for different coffee preferences;
    \item \(C_2\): \textit{coffee beans by roast and flavor profile}, covering roast level, tasting notes, origin, acidity, body, and recommended brewing style;
    \item \(C_3\): \textit{coffee machine compatibility and brewing guidance}, covering which coffees fit which machines or brewing methods, such as espresso machines, filter brewing, pods, or automatic machines;
    \item \(C_4\): \textit{coffee quiz and guided product selection}, covering pages that help users identify coffee preferences through question-based guidance;
    \item \(C_5\): \textit{bundles, gifts, and curated offers}, covering bundles, seasonal packs, gift cards, and one-time purchase flows;
    \item \(C_6\): \textit{coffee accessories and related products}, covering machines, thermoses, accessories, and related shopping journeys.
\end{itemize}

For \(C=\) ``coffee subscriptions and recurring delivery,'' \(P_C^{\text{brand}}\) may include pages about subscription plans, delivery cadence, product selection, and relevant promotions, while \(R_C\) contains comparable external pages on the same topic.

For a topic cluster \(C\), let
\[
P_C^{\text{brand}}=\{p_1^{\text{brand}},p_2^{\text{brand}},\dots,p_k^{\text{brand}}\}
\]
be the set of brand-controlled page snapshots assigned to the cluster, and let
\[
R_C = \{p_1^{\mathrm{ref}}, p_2^{\mathrm{ref}}, \dots, p_n^{\mathrm{ref}}\}
\]
be a reference set of comparable external pages on the same topic.

Information Gain is computed in two stages. Each page snapshot first receives page-level dimension estimates, optionally compared with \(R_C\), normalized to \([0,1]\). These values are then aggregated by topic cluster to produce cluster-level coverage, consistency, and readiness indicators. Recommendations are indexed by cluster and linked to affected pages.

For a page snapshot \(p\), let \(\hat{f}_j(p,C,R_C)\) denote the page-level estimate for dimension \(j\). The page-level score is
\[
\hat{g}(p \mid C,R_C)
=
\sum_j w^{\text{page}}_{C,j}\hat{f}_j(p,C,R_C),
\qquad
\sum_j w^{\text{page}}_{C,j}=1,
\qquad
w^{\text{page}}_{C,j}\ge 0.
\]
The corresponding cluster-level feature value is obtained through a dimension-specific aggregation rule:
\[
\hat{F}_j(C,R_C)
=
\mathcal{A}_j
\left(
\{\hat{f}_j(p,C,R_C):p\in P_C^{\text{brand}}\}
\right).
\]
The aggregation rule may use a weighted mean, a maximum over coverage signals, a union of claims, a cross-page consistency check, or a penalty for duplicated or risky content, depending on the dimension.

Operational page and cluster scores use the same named dimensions, with lowercase symbols for page estimates and uppercase symbols for cluster aggregates:
\[
\mathbf{F}_{\text{v1}}(C,R_C)
=
\langle
F_{\text{intent}},
F_{\text{spec}},
F_{\text{str}},
F_{\text{trust}},
F_{\text{orig}},
F_{\text{search}},
F_{\text{media}},
F_{\text{risk}}
\rangle.
\]

The dimensions are interpreted as follows:
\begin{itemize}
    \item \(\mathbf{F_{\text{intent}}}\) --- \textbf{intent coverage.} This dimension is high when the cluster collectively satisfies the main user needs associated with the topic. Typical indicators include coverage of definitions, product selection criteria, pricing or availability expectations, constraints, policies, comparisons, and follow-up questions likely to arise for the same intent.
    \item \(\mathbf{F_{\text{spec}}}\) --- \textbf{specificity and boundedness.} This dimension is high when the cluster states concrete facts rather than generic claims. Typical indicators include explicit numbers, prices, delivery timing, quantities, named products, roast levels, flavor notes, compatibility constraints, applicability conditions, exclusions, limits, and failure cases.
    \item \(\mathbf{F_{\text{str}}}\) --- \textbf{structured answerability.} This dimension is high when the cluster exposes information in reusable units, such as headings, concise sections, ordered steps, bullet lists, comparison tables, checklists, product cards, and FAQ entries.
    \item \(\mathbf{F_{\text{trust}}}\) --- \textbf{evidence and trust support.} This dimension is high when claims are supported by product facts, examples, source references, policy details, contact or organizational signals, certifications, measurements, observed outcomes, or other provenance cues, and when claims are internally consistent across the cluster.
    \item \(\mathbf{F_{\text{orig}}}\) --- \textbf{originality and information gain relative to alternatives.} This dimension is high when the cluster contributes useful claims, examples, comparisons, or structured details not already common in comparable pages in \(R_C\). It combines differentiation and novelty into one operational dimension for V1.
    \item \(\mathbf{F_{\text{search}}}\) --- \textbf{search and discovery readiness.} This dimension is high when the content can be discovered, rendered, indexed, and interpreted by general search and AI-mediated discovery systems. Typical indicators include crawlable and renderable content, descriptive titles and metadata, clear canonical signals, internal links, stable URLs, non-duplicative page organization, and content visible in the rendered page.
    \item \(\mathbf{F_{\text{media}}}\) --- \textbf{media and product readiness.} This dimension is high when images, product names, product descriptions, prices, availability cues, and product-specific attributes are clear, close to the relevant text, and consistent across listing and detail pages. Structured product data may be added as an implementation extension.
    \item \(\mathbf{F_{\text{risk}}}\) --- \textbf{content risk control.} This dimension is high when the cluster avoids low-value or search-engine-first patterns, such as thin mass-produced pages, excessive duplication, keyword stuffing, artificial freshness, unsupported generated text, contradictory claims, and pages created only to capture small query variations without adding user value.
\end{itemize}

The cluster-level score is computed as a weighted combination of the aggregated dimensions:
\[
\hat{G}(C \mid R_C)=\sum_j w_{C,j}\hat{F}_j(C,R_C),
\qquad
\sum_j w_{C,j} = 1,
\qquad
w_{C,j}\ge 0.
\]
The page weights \(w^{\text{page}}_{C,j}\) and cluster weights \(w_{C,j}\) are cluster-specific configuration values stored in the active Information Gain Card.

Each dimension may be estimated through deterministic heuristics, human labels, LLM-based labels, or a hybrid process. The first implementation uses heuristic extraction, with hybrid judging reserved for dimensions requiring semantic calibration, especially intent coverage, trust support, and originality relative to a reference set.

\begin{figure}[H]
    \centering
    \includegraphics[width=0.49\linewidth]{resources/content_sidescoring.png}
    \caption{Content-side static scoring pipeline. Brand pages are captured as versioned snapshots, assigned to topic clusters, and tagged with the active variant when applicable. Feature extraction computes page-level scores, cluster aggregation combines page evidence into topic-level indicators, and both levels are evaluated against static rules stored in the Information Gain Card. Failing rules are routed to the recommendation workflow with affected pages retained as edit targets.}
    \label{fig:static-scoring-pipeline}
\end{figure}

\paragraph{Weight optimization.}
In the first implementation phase, page weights \(w^{\text{page}}_{C,j}\) and cluster weights \(w_{C,j}\) are set manually and kept fixed for interpretability. In the second phase, once sufficient page scores, cluster scores, and runtime QAT data have been collected, the weights may be calibrated against downstream metrics such as alignment, presence, citation share, or observed site outcomes, subject to the constraints \(w^{\text{page}}_{C,j} \ge 0\), \(w_{C,j} \ge 0\), and both weight vectors summing to one. In the third phase, weights may be specialized by cluster.

\subsection{Information Gain Card (IGC)}

For each topic cluster \(C\), the project maintains an Information Gain Card
\[
IGC_C = \langle \Sigma_C^{\text{page}}, \Sigma_C^{\text{cluster}}, \Phi_C, \tau_C, \Gamma_C, \Omega_C \rangle.
\]

The components are tied to different stages of evaluation and intervention:
\begin{itemize}
    \item \(\Sigma_C^{\text{page}}\) --- \textbf{page-level static requirements}, evaluated on each page snapshot \(p \in P_C^{\text{brand}}\) at authoring time, publish time, or immediately after revision. These rules are configuration-driven and may include:
    \[
    \hat{g}(p \mid C,R_C) < \theta_C^{\text{page}},
    \qquad
    \hat{f}_{\text{spec}}(p,C) < \theta_C^{\text{page-spec}},
    \qquad
    \hat{f}_{\text{str}}(p,C) < \theta_C^{\text{page-str}},
    \qquad
    \hat{f}_{\text{search}}(p,C) < \theta_C^{\text{page-search}}.
    \]
    A page failing any active rule is marked as an affected page for the cluster recommendation workflow.

    \item \(\Sigma_C^{\text{cluster}}\) --- \textbf{cluster-level static requirements}, evaluated on the aggregated cluster evidence set \(P_C^{\text{brand}}\). These rules are configuration-driven and may include:
    \[
    \hat{G}(C \mid R_C) < \theta_C^{\text{IG}},
    \qquad
    \hat{F}_{\text{intent}}(C) < \theta_C^{\text{intent}},
    \qquad
    \hat{F}_{\text{spec}}(C) < \theta_C^{\text{spec}},
    \qquad
    \hat{F}_{\text{str}}(C) < \theta_C^{\text{str}},
    \qquad
    \hat{F}_{\text{trust}}(C) < \theta_C^{\text{trust}},
    \qquad
    \hat{F}_{\text{search}}(C) < \theta_C^{\text{search}}.
    \]
    A cluster failing any active rule is marked \texttt{NeedsRevision} and routed to the content recommendation workflow. The recommendation is cluster-indexed, while page-level rule failures identify the concrete surfaces to edit.
 
    \item \(\Phi_C\) --- \textbf{runtime expectations}, evaluated on aggregated QAT metrics after answer collection. This layer measures whether AI systems expose the intended behavior for the cluster through explicit target thresholds for presence, prominence, citation share, and alignment.

    \item \(\tau_C\) --- \textbf{temporal review rules}, evaluated across consecutive measurement windows. This layer determines when a runtime weakness is persistent enough to justify a new intervention cycle after publication or revision.

    \item \(\Gamma_C\) --- \textbf{canonical claims}, used to evaluate whether answers and content surfaces express the cluster's intended narrative, constraints, product facts, and policy claims.

    \item \(\Omega_C\) --- \textbf{recommendation mappings}, used to translate failed static dimensions, runtime failures, and risk signals into cluster-specific content actions.
\end{itemize}

Operationally, \(\Sigma_C^{\text{page}}\) runs on individual page snapshots, \(\Sigma_C^{\text{cluster}}\) runs on the aggregated cluster evidence set, \(\Phi_C\) runs on aggregated answer traces, and \(\tau_C\) runs on sequences of aggregated windows. In implementation terms, \(IGC_C\) acts as a cluster-specific policy object. It stores page weights, cluster weights, thresholds, review rules, canonical claims, page-type expectations, and recommendation mappings used by the evaluation pipeline. Different clusters can therefore apply different standards, for example stricter evidence requirements for product-comparison content, stricter structured-answerability requirements for guided-selection or FAQ-oriented content, and stricter media or product-readiness requirements for commerce-heavy clusters.

\subsection{Query--Answer Traces (QAT)}

A Query--Answer Trace records one executed query and the corresponding AI answer:
\[
r = \langle t, q, s, a, D \rangle,
\]
where \(t\) is the timestamp, \(q\) is the query, \(s\) is the AI system, \(a\) is the answer text, and \(D\) is the set of cited domains or URLs.

Each trace is evaluated through four core runtime metrics:
\begin{itemize}
    \item \textbf{presence}, aggregated as \(\operatorname{PresenceRate}(C,T)\);
    \item \textbf{prominence}, aggregated as \(\operatorname{ProminenceMean}(C,T)\);
    \item \textbf{citation share}, aggregated as \(\operatorname{CitationShareMean}(C,T)\);
    \item \textbf{alignment}, aggregated as \(\operatorname{AlignmentMean}(C,T)\).
\end{itemize}

Let \(B\) denote the brand identifier set associated with the cluster, including the official brand domain and the accepted textual brand name. Let
\[
\Gamma_C=\{\gamma_1,\dots,\gamma_m\}
\]
denote the canonical claim set stored in \(IGC_C\), that is, the set of brand-relevant claims considered correct and important for answers in that cluster.

The trace-level metrics are defined by
\begin{align}
m_{\text{pres}}(r)
&=
\mathbf{1}\{D \cap B \neq \emptyset \;\lor\; \text{an element of } B \text{ appears in } a\},\\[4pt]
m_{\text{prom}}(r)
&=
\begin{cases}
1/\operatorname{rank}_B(r), & \text{if the brand appears},\\
0, & \text{otherwise},
\end{cases}\\[4pt]
m_{\text{cite}}(r)
&=
\frac{|D \cap B|}{|D|+\varepsilon},\\[4pt]
m_{\text{align}}(r)
&=
\frac{1}{|\Gamma_C|}
\sum_{i=1}^{|\Gamma_C|}
\operatorname{match}(\gamma_i,a).
\end{align}

If \(R_{C,T}\) denotes the set of traces collected for cluster \(C\) in measurement window \(T\), the aggregated runtime indicators are
\begin{align}
\operatorname{PresenceRate}(C,T)
&=
\frac{1}{|R_{C,T}|}\sum_{r\in R_{C,T}} m_{\text{pres}}(r),\\[4pt]
\operatorname{ProminenceMean}(C,T)
&=
\frac{1}{|R_{C,T}|}\sum_{r\in R_{C,T}} m_{\text{prom}}(r),\\[4pt]
\operatorname{CitationShareMean}(C,T)
&=
\frac{1}{|R_{C,T}|}\sum_{r\in R_{C,T}} m_{\text{cite}}(r),\\[4pt]
\operatorname{AlignmentMean}(C,T)
&=
\frac{1}{|R_{C,T}|}\sum_{r\in R_{C,T}} m_{\text{align}}(r).
\end{align}

Sentiment may be stored as an auxiliary signal, but it is not part of the core implementation.

\subsection{Failure Taxonomy for IG in AI-Mediated Discovery}

The failure taxonomy is used to diagnose why a cluster underperforms and to decide which type of intervention is needed: content revision, runtime re-evaluation, or cross-channel consistency review.

\begin{itemize}
    \item \textbf{IG-F1: Visibility Gap.}
    The cluster scores well under the static content criteria, but runtime visibility remains weak:
    \[
    \hat{G}(C \mid R_C)\ge\theta_C^{\text{IG}}
    \quad \land \quad
    \big(
    \operatorname{PresenceRate}(C,T)<\theta_C^{\text{pres}}
    \;\lor\;
    \operatorname{ProminenceMean}(C,T)<\theta_C^{\text{prom}}
    \big).
    \]

    \item \textbf{IG-F2: Narrative Misalignment.}
    The brand appears in answers, but the answer content does not match the canonical claims for the cluster:
    \[
    \operatorname{PresenceRate}(C,T)\ge\theta_C^{\text{pres}}
    \quad \land \quad
    \operatorname{AlignmentMean}(C,T)<\theta_C^{\text{align}}.
    \]

    \item \textbf{IG-F3: Hollow IG.}
    The page set appears superficially rich, for example through length, repetition, or many low-value surfaces, but remains weak in specificity, evidence, structure, originality, or risk control:
    \[
    \hat{G}(C \mid R_C)<\theta_C^{\text{IG}}
    \quad \land \quad
    \big(
    \operatorname{Length}(P_C^{\text{brand}})>\theta^{\text{len}}
    \;\lor\;
    \operatorname{KeywordFrequency}(P_C^{\text{brand}})>\theta^{\text{kw}}
    \;\lor\;
    \operatorname{Duplication}(P_C^{\text{brand}})>\theta^{\text{dup}}
    \big).
    \]

    \item \textbf{IG-F4: IG Drift after Revision.}
    A page or cluster revision reduces the corresponding IG score and is followed by degraded runtime behavior:
    \[
    \left(
    \hat{G}(C^{\text{post}} \mid R_C) - \hat{G}(C^{\text{pre}} \mid R_C) < 0
    \;\lor\;
    \exists p \in P_C^{\text{brand}}:
    \hat{g}(p^{\text{post}} \mid C,R_C)-\hat{g}(p^{\text{pre}} \mid C,R_C)<0
    \right)
    \quad \land \quad
    \big(
    \Delta \operatorname{AlignmentMean}(C,T) < 0
    \;\lor\;
    \Delta \operatorname{PresenceRate}(C,T) < 0
    \big).
    \]

    \item \textbf{IG-F5: Cross-Channel Inconsistency.}
    Brand-controlled pages assigned to the same cluster contain conflicting claims, values, or capability descriptions, leading to unstable or contradictory answer narratives across traces.

    \item \textbf{IG-F6: Discovery Readiness Failure.}
    The cluster may contain useful content, but one or more page-level or cluster-level discoverability or interpretation conditions are weak, such as crawlability, rendered visibility, canonical clarity, internal linking, media context, product-data consistency, or duplicate control:
    \[
    \hat{F}_{\text{search}}(C)<\theta_C^{\text{search}}
    \quad \lor \quad
    \hat{F}_{\text{media}}(C)<\theta_C^{\text{media}}
    \quad \lor \quad
    \exists p \in P_C^{\text{brand}}:
    \hat{f}_{\text{search}}(p,C)<\theta_C^{\text{page-search}}.
    \]
\end{itemize}

\section{Methodology}

\subsection{Content Sampling and IG-Boosting Interventions}

For each brand or product area, the study defines a set of topic clusters
\[
\mathcal{C}=\{C_1,C_2,\dots,C_K\},
\]
typically with \(K\) between \(5\) and \(10\). Each cluster \(C\) contains:
\begin{itemize}
    \item a set of brand-controlled pages;
    \item a reference set \(R_C\) of comparable external pages;
    \item a cluster-specific Information Gain Card \(IGC_C\).
\end{itemize}

Within each cluster, a subset of brand-controlled pages is selected for intervention. The content unit of change is the page revision, because authors and product teams edit concrete URLs, blocks, product records, and page metadata. The evaluation then measures both the edited page and the containing cluster. For a treated cluster \(C\), a revised evidence set \(P_C^{\text{post}}\) is produced with the objective of increasing the relevant page-level scores and the cluster-level Information Gain score,
\[
\hat{g}(p^{\text{post}} \mid C,R_C) > \hat{g}(p^{\text{pre}} \mid C,R_C)
\quad\text{for edited pages }p,
\]
and
\[
\hat{G}(C^{\text{post}} \mid R_C) > \hat{G}(C^{\text{pre}} \mid R_C),
\]
while preserving the original topic intent.

The revision targets the dimensions introduced above, especially intent coverage, specificity, structured answerability, evidence and trust support, originality, search readiness, media and product readiness, and content risk control. Typical interventions include:
\begin{itemize}
    \item adding concrete examples or explicit use cases;
    \item adding numeric claims, bounded statements, or measurable outcomes;
    \item adding explicit applicability conditions, exclusions, or cases in which the recommendation does not hold;
    \item restructuring the page into lists, steps, tables, checklists, or FAQ sections;
    \item adding concise differentiators relative to competing solutions;
    \item clarifying titles, metadata, canonical page roles, and internal links;
    \item improving product names, descriptions, image context, availability cues, and price consistency;
    \item removing duplicated, thin, unsupported, or search-engine-first content patterns.
\end{itemize}

In the intended product setting, the recommendation engine operates over cluster-indexed page snapshots: it detects low-scoring pages and low-scoring cluster dimensions, maps them to recommended actions, and emits revision guidance tied to the affected pages. The evaluation then compares page versions for local content quality and cluster versions for aggregate IG scores and QAT-derived runtime metrics.

\subsection{IG Rubric and IGC Instantiation}

Each captured page is scored on a rubric whose dimensions correspond to the IG features defined above: intent coverage, specificity, structured answerability, evidence and trust support, originality, search readiness, media and product readiness, and content risk control. The cluster score aggregates those page-level rubric values and adds cross-page signals such as coverage completeness, claim consistency, duplication, and page-role clarity. Judges or automated heuristics may assign discrete scores, for example on a \(0\)--\(4\) scale, and the normalized feature values are obtained by dividing by the maximum score.

The operational IG score is computed as
\[
\hat{g}(p \mid C,R_C)=\sum_j w^{\text{page}}_{C,j}\hat{f}_j(p,C),
\qquad
\sum_j w^{\text{page}}_{C,j} = 1,
\]
and
\[
\hat{G}(C \mid R_C)=\sum_j w_{C,j}\hat{F}_j(C),
\qquad
\sum_j w_{C,j} = 1.
\]

Weight selection follows the calibration procedure described in Section~\ref{sec:igscores}.

For each cluster, the Information Gain Card is instantiated by specifying:
\begin{itemize}
    \item page-level static requirements for editable/indexable surfaces;
    \item cluster-level static requirements on the aggregated evidence set;
    \item runtime targets for visibility and alignment;
    \item temporal review rules for opening a new intervention cycle;
    \item canonical claims and page-type expectations;
    \item cluster-specific recommendation mappings.
\end{itemize}

\subsection{Recommendation Output Layer}

For each evaluated cluster \(C\), the system generates a cluster-specific recommendation bundle
\[
\mathcal{R}(C)=\langle \text{issues},\text{evidence},\text{actions},\text{affected pages},\text{priority}\rangle,
\]
where:
\begin{itemize}
    \item \textbf{issues} identifies failed dimensions or cluster rules, such as low intent coverage, weak specificity, low trust support, low search readiness, or low runtime alignment;
    \item \textbf{evidence} contains the cluster score breakdown, affected page snapshots, relevant trace examples, and failed thresholds;
    \item \textbf{actions} contains recommended content changes, such as adding a comparison table, clarifying product applicability, introducing explicit quantitative details, resolving conflicting claims, improving internal links, or clarifying product media and descriptions;
    \item \textbf{affected pages} identifies the pages or content surfaces where the action should be applied;
    \item \textbf{priority} captures implementation urgency based on page importance, cluster importance, and runtime failure severity.
\end{itemize}

Recommendation generation is initially rule-based and derives actions from cluster-specific recommendation mappings \(\Omega_C\), low-scoring page or cluster dimensions, and failure categories. For example, low \(f_{\text{spec}}\) on an affected page or low \(F_{\text{spec}}\) across a cluster may trigger recommendations to add concrete constraints, numerical values, or explicit eligibility conditions on named pages; low structure may trigger recommendations to add headings, lists, tables, or FAQ blocks; low search readiness may trigger recommendations to clarify titles, canonical roles, internal links, or duplicated page purposes; low media readiness may trigger recommendations to improve product descriptions, image context, or price consistency; and IG-F2 may trigger recommendations to rewrite affected pages so that canonical claims become easier for AI systems to recover accurately.

\begin{table}[H]
\centering
\small
\begin{tabular}{@{}p{0.27\linewidth} p{0.34\linewidth} p{0.30\linewidth}@{}}
\textbf{Surface} & \textbf{Change type} & \textbf{Primary role} \\
\hline
AEM-authored content update & Canonical page copy, headings, blocks, FAQs, links, and metadata. & Durable content revision. \\
CMS workflow & Review, approval, publishing, ownership, and audit trail. & Governance of content changes. \\
Edge/CDN variant & Controlled alternate response or transformation served by routing rules. & Variant testing or targeted exposure before canonical revision. \\
\end{tabular}
\caption{Intervention surfaces for recommendation deployment.}
\label{tab:intervention-surfaces}
\end{table}

\begin{table}[H]
\centering
\small
\begin{tabular}{@{}p{0.24\linewidth} p{0.36\linewidth} p{0.31\linewidth}@{}}
\textbf{Failed signal} & \textbf{Recommendation action} & \textbf{Intervention surface} \\
\hline
Low specificity & Add prices, timing, quantities, constraints, product facts, or eligibility conditions. & AEM copy or product data. \\
Low structure & Add headings, lists, comparison tables, steps, or FAQ-like sections. & AEM-authored page blocks. \\
Low media readiness & Improve product name, description, price, image context, or PLP/PDP consistency. & Product or commerce data. \\
Low search readiness & Clarify title, metadata, canonical role, internal links, or duplicate page purpose. & AEM metadata, page structure, or routing configuration. \\
IG-F2 narrative misalignment & Rewrite affected pages around canonical claims. & AEM copy, PDP content, or reviewed CMS update. \\
\end{tabular}
\caption{Mapping from failed signals to actions and intervention surfaces.}
\label{tab:recommendation-intervention-map}
\end{table}

\begin{table}[H]
\centering
\small
\begin{tabular}{@{}p{0.22\linewidth} p{0.26\linewidth} p{0.25\linewidth} p{0.18\linewidth}@{}}
\textbf{Diagnosis} & \textbf{Recommendation} & \textbf{Intervention} & \textbf{Evaluation} \\
\hline
Low specificity on \texttt{/coffee}. & Add roast level, package size, flavor notes, and brewing guidance. & Update product data or authored copy. & Compare page and cluster score deltas. \\
\end{tabular}
\caption{Example recommendation-to-intervention trace.}
\label{tab:recommendation-example}
\end{table}

Later versions may add LLM-assisted rewriting, agentic revision workflows, or learned prioritization of candidate interventions.

\begin{figure}[H]
    \centering
    \includegraphics[width=0.49\linewidth]{resources/recomsystem.png}
    \caption{Recommendation and deployment loop. Static failures from page-level scoring, cluster aggregation, and runtime AI-answer evaluation are combined in the recommendation engine, which produces ranked, cluster-specific content actions with affected pages for review. After reviewer approval, recommendations can be deployed either through origin-side CMS updates or through controlled edge/CDN variants for agent-facing traffic. The deployed changes then enter the next measurement cycle, where pages and clusters are re-scored and runtime behavior is re-evaluated. This closes the loop between diagnosis, intervention, deployment, and measurement.}
    \label{fig:recommendation-deployment-loop}
\end{figure}

\subsection{Human and Hybrid LLM-Based Annotation Strategy}

Annotation is performed on two object types:
\begin{itemize}
    \item \textbf{page evidence objects} \((p,C,R_C)\), used to estimate page-level IG dimensions and the derived score \(\hat{g}(p \mid C,R_C)\);
    \item \textbf{cluster evidence objects} \((C,P_C^{\text{brand}},R_C)\), used to aggregate page labels, estimate cross-page dimensions, and derive the cluster score \(\hat{G}(C \mid R_C)\);
    \item \textbf{trace objects} \(r\), used to label runtime properties of QAT records, especially alignment with \(\Gamma_C\) and the failure categories IG-F1--IG-F6.
\end{itemize}

For a page evidence object, the annotation output is
\[
y^{\text{page}}(p)
=
\langle
y_{\text{intent}},
y_{\text{spec}},
y_{\text{str}},
y_{\text{trust}},
y_{\text{orig}},
y_{\text{search}},
y_{\text{media}},
y_{\text{risk}}
\rangle,
\]
where each component is a discrete rubric score. The page label is used directly for local edit targeting. For a cluster evidence object, the annotation output is
\[
y^{\text{cluster}}(C)
=
\langle
y_{\text{intent}},
y_{\text{spec}},
y_{\text{str}},
y_{\text{trust}},
y_{\text{orig}},
y_{\text{search}},
y_{\text{media}},
y_{\text{risk}}
\rangle,
\]
where each component is the aggregated cluster label, optionally including cross-page consistency, duplication, coverage, and page-role signals.

For a trace object \(r\), the annotation output is
\[
y^{\text{trace}}(r)
=
\langle
y_{\text{align}},
y_{\text{failure}}
\rangle,
\]
where \(y_{\text{align}}\) is the alignment judgment with respect to \(\Gamma_C\), and \(y_{\text{failure}}\) is a set of zero or more tags from \(\{\text{IG-F1},\dots,\text{IG-F6}\}\).

The judge pool is
\[
J = J_{\text{human}} \cup J_{\text{llm}}.
\]
Each judge \(j \in J\) emits a label \(\hat{y}^{(j)}(z)\) for an object \(z \in \{p,C,r\}\). The final label is obtained by aggregation:
\[
\bar{y}(z)=\mathcal{A}\big(\{\hat{y}^{(j)}(z)\}_{j \in J}\big),
\]
where the aggregation rule depends on label type:
\begin{itemize}
    \item ordinal rubric scores are combined with a median or robust mean;
    \item categorical labels are combined with majority vote;
    \item failure tags are combined with a confidence-thresholded union.
\end{itemize}

The workflow is staged as follows:
\begin{enumerate}
    \item \textbf{Human reference phase.} A domain specialist or content expert labels a gold subset of page evidence objects, cluster evidence objects, and traces manually.
    \item \textbf{Calibration phase.} LLM-based judges are run on the same gold subset and compared against the human labels.
    \item \textbf{Operational phase.} Most page and trace objects are labeled by calibrated LLM judges, while humans review sampled items, disagreement cases, and high-priority failures.
\end{enumerate}

An object \(z\) is escalated to human review if at least one of the following conditions holds:
\[
\operatorname{Disagreement}(z) > \lambda,
\qquad
\operatorname{Confidence}(z) < \mu,
\qquad
\operatorname{Priority}(z) \in \{\texttt{P0},\texttt{P1}\}.
\]
In V1, \(\operatorname{Priority}(z)\) is configuration-driven and may be set to \texttt{P0} or \texttt{P1} for homepage content, cluster landing pages, pages attached to active experiments, or traces that produce severe brand inaccuracies.

\subsection{Experimental Design: Before--After Evaluation}

The primary evaluation design is a before--after comparison that preserves the page as the intervention unit and the cluster as the aggregation unit. Let
\[
P_C^{\text{pre}}, P_C^{\text{post}} \subseteq P_C^{\text{brand}}
\]
denote the pre-intervention and post-intervention evidence sets for cluster \(C\), and let \(p^{\text{pre}}\) and \(p^{\text{post}}\) denote an edited page before and after intervention. The study compares pre-intervention and post-intervention values of the outcome variable \(Y\), where \(Y\) may denote:
\begin{itemize}
    \item a \emph{page static content outcome}, such as \(\hat{g}(p \mid C,R_C)\) or one of the page dimensions \(\hat{f}_j(p,C)\);
    \item a \emph{cluster static content outcome}, such as \(\hat{G}(C \mid R_C)\) or one of the cluster dimensions \(\hat{F}_j(C)\);
    \item a \emph{visibility outcome}, such as presence or prominence;
    \item a \emph{narrative outcome}, such as alignment;
    \item a \emph{site outcome}, such as agentic referrals, visits, engagement, or tracked business actions, provided that the exposed content variant can be joined to delivery logs and downstream event data.
\end{itemize}

For a treated cluster \(C\), the before--after change is
\[
\Delta^{\text{BA}}_C
=
\bar{Y}^{\text{post}}_C-\bar{Y}^{\text{pre}}_C.
\]
For an edited page \(p\), the local before--after change is
\[
\Delta^{\text{BA}}_p
=
\bar{Y}^{\text{post}}_p-\bar{Y}^{\text{pre}}_p.
\]
The page-level change measures the edited surface. The cluster-level change measures the aggregate topic outcome across page and runtime evidence.

V1 uses before--after comparison as the primary design, while preserving sufficient metadata for later matched-control or quasi-experimental refinement.

\subsection{Synthetic Prompts and QAT Collection}

In V1, the existing prompt inventory may be used as the initial prompt set and mapped to the broad \texttt{unknown} cluster. As more clusters are introduced, each prompt is assigned to the most specific cluster it tests, and historical traces can be reprocessed when the mapping changes.

For each topic cluster \(C\), the study defines a prompt set \(Q_C\) covering recurring user intents such as:
\begin{itemize}
    \item \textbf{definition}, for example: ``What is a coffee subscription?'', ``What does medium roast mean?'', or ``What is the purpose of a coffee quiz?'';
    \item \textbf{comparison}, for example: ``How is a subscription different from a one-time coffee purchase?'', ``Which is better for espresso, a dark roast or a medium roast?'', or ``How does a bundle compare with buying items separately?'';
    \item \textbf{product selection}, for example: ``Which coffee is best for an automatic machine?'', ``Which subscription is best for someone who likes espresso?'', or ``What gift option should I choose for a coffee lover?'';
    \item \textbf{how-to guidance}, for example: ``How do I choose the right coffee for filter brewing?'', ``How can I pause or modify a subscription?'', or ``How do I decide between pods, beans, and ground coffee?''
\end{itemize}

For Frescopa-like deployments, these prompts can be specialized to observable surfaces such as subscription selection, roast guidance, machine compatibility, coffee quizzes, bundles, gift-oriented offers, and accessories.

For each cluster \(C\), the prompt set \(Q_C\) is executed against one or more AI systems across repeated measurement windows. For each query--system pair, the system stores one or more answer samples as QAT records. When feasible, prompts are repeated within the same window to reduce single-run variance. The resulting traces are grouped by cluster and time window, then aggregated into the runtime indicators used by \(\Phi_C\) and \(\tau_C\).

\begin{algorithm}[H]
\caption{QAT collection for a topic cluster}
\KwIn{Cluster \(C\), prompt set \(Q_C\), AI systems \(S\), measurement windows \(W\), repetition count \(n\)}
\KwOut{QAT records \(R_{C,W}\) and runtime indicator tuples \(M_{C,w}\)}

Initialize \(R_{C,W} \leftarrow \emptyset\)

\ForEach{measurement window \(w \in W\)}{
    \ForEach{prompt \(q \in Q_C\)}{
        \ForEach{AI system \(s \in S\)}{
            \For{\(i \leftarrow 1\) \KwTo \(n\)}{
                Execute \(q\) on \(s\)\;
                Store answer as QAT record \(r=\langle t,q,s,a,D\rangle\)\;
                Add \(r\) to \(R_{C,w}\)\;
            }
        }
    }
    Compute the runtime indicator tuple
    \[
    M_{C,w}
    =
    \langle
    \operatorname{PresenceRate}(C,w),
    \operatorname{ProminenceMean}(C,w),
    \operatorname{CitationShareMean}(C,w),
    \operatorname{AlignmentMean}(C,w)
    \rangle
    \]
    from \(R_{C,w}\)\;
}
\end{algorithm}

\begin{figure}[H]
    \centering
    \includegraphics[width=0.49\linewidth]{resources/runtime_aianswers.png}
    \caption{Runtime AI-answer evaluation pipeline. A versioned prompt inventory is scheduled across measurement windows and executed against selected AI systems. The resulting answers are stored as Query--Answer Traces, which are then processed by the judging and calibration layer to assign alignment and failure labels. Aggregated runtime metrics, such as presence, prominence, citation share, and alignment, are evaluated against runtime and temporal rules defined in the cluster card. Healthy clusters remain in monitoring mode, while persistent failures open a new intervention cycle.}
    \label{fig:runtime-evaluation-pipeline}
\end{figure}

\subsection{Phase 2 Optional Extension: Adversarial Prompt Exploration}
This module is out of scope for the first implementation cycle and retained as a Phase 2 extension. When enabled, it starts from the baseline prompt set \(Q_C\) for cluster \(C\) and generates additional prompts intended to reveal weaknesses not exposed by the baseline inventory.

At each iteration \(t\), the system observes the current search state
\[
h_t=\langle C,\;Q_C,\;q_1,\dots,q_{t-1},\;r_1,\dots,r_{t-1}\rangle,
\]
which includes the active cluster, the baseline prompts, the prompts already tested, and their associated QAT records.

The action at step \(t\) is a prompt transformation selected from a small set of operators, for example:
\begin{itemize}
    \item rewriting an existing prompt;
    \item changing the user-intent framing;
    \item adding or removing a constraint;
    \item changing the usage scenario while remaining within the same cluster.
\end{itemize}

The chosen action produces a candidate prompt \(q_t\), which is executed against the target AI system and stored as a new QAT record.

Candidate prompts are scored by:
\begin{itemize}
    \item a \textbf{reward} \(R_t\), which increases when the resulting answer exposes a weakness, such as low brand presence, low alignment, or dominant competitor visibility;
    \item a \textbf{cost} \(C_t\), which penalizes unrealistic, excessively long, or off-cluster prompts.
\end{itemize}

The optimization objective is
\[
\mathbb{E}\left[\sum_{t=1}^{T}\gamma^t(R_t-\beta C_t)\right],
\]
where \(T\) is the exploration horizon, \(\gamma\) is the discount factor, and \(\beta\) controls the penalty applied to low-quality prompts.

\subsection{Toy Example}

Consider a brand page in the topic cluster \(C=\) ``coffee subscriptions and recurring delivery.''

\subsection*{Baseline version}

\emph{``Our subscription helps coffee lovers receive fresh coffee regularly. Customers can choose products they like and enjoy convenient delivery.''}

This version is generic. It does not specify delivery cadence, pause conditions, product-selection criteria, evidence for freshness, or clear differences between plans. Under the rubric, the edited page would receive low scores on \(f_{\text{intent}}\), \(f_{\text{spec}}\), \(f_{\text{str}}\), and \(f_{\text{trust}}\). If similar gaps appear across related pages, the cluster would also receive low scores on \(F_{\text{intent}}\), \(F_{\text{spec}}\), \(F_{\text{str}}\), and \(F_{\text{trust}}\).

A cluster card may require both page-level and cluster-level static rules:
\[
\Sigma_C^{\text{page}}:
\hat{g}(p \mid C,R_C)\ge\theta_C^{\text{page}},
\quad
\hat{f}_{\text{spec}}(p,C)\ge\theta_C^{\text{page-spec}},
\quad
\hat{f}_{\text{str}}(p,C)\ge\theta_C^{\text{page-str}};
\]
\[
\Sigma_C^{\text{cluster}}:
\hat{G}(C \mid R_C)\ge\theta_C^{\text{IG}},
\quad
\hat{F}_{\text{intent}}(C)\ge\theta_C^{\text{intent}},
\quad
\hat{F}_{\text{spec}}(C)\ge\theta_C^{\text{spec}},
\quad
\hat{F}_{\text{str}}(C)\ge\theta_C^{\text{str}};
\]
\[
\Phi_C:
\operatorname{PresenceRate}(C,T) \ge \theta_C^{\text{pres}}
\quad\text{and}\quad
\operatorname{AlignmentMean}(C,T) \ge \theta_C^{\text{align}};
\]
\[
\tau_C:
\text{trigger review if either runtime metric remains below threshold for }k\text{ consecutive windows.}
\]

\subsection*{IG-boosted version}

\emph{``Our coffee subscription delivers freshly roasted coffee every two or four weeks. Customers can pause or modify the plan after the second delivery. The subscription includes guidance for selecting coffees based on roast preference and brewing style. Each plan page includes roast level, tasting notes, and recommended usage, together with a comparison table for 250g, 500g, and 1kg options.''}

This version introduces explicit delivery rules, applicability conditions, product-selection guidance, and reusable structure. The expected result is an increase in the edited page score \(\hat{g}(p \mid C,R_C)\) and, when the information is consistent across related pages, an increase in the cluster-level score \(\hat{G}(C \mid R_C)\).

\section{Evaluation Plan}

Evaluation is carried out at three layers:
\begin{itemize}
    \item \textbf{page static level}, through changes in page IG score, dimension breakdowns, affected-page status, and static compliance before and after intervention;
    \item \textbf{cluster static level}, through changes in aggregated cluster IG score, dimension breakdowns, recommendation status, and static compliance before and after intervention;
    \item \textbf{cluster runtime level}, through changes in QAT indicators such as presence, prominence, citation share, and alignment.
\end{itemize}

When available, the evaluation also includes site outcomes such as agentic referrals, visits, engagement, and tracked conversion actions (for example subscription starts, completed purchases, newsletter sign-ups, quiz completions, or add-to-cart events). The main sources of uncertainty are variation across AI systems, delayed or unstable indexing behavior, imperfect traffic attribution, and residual disagreement between human and model-based judges.

\subsection{Site Outcome Attribution and Agentic Referral Measurement}

In V1, the system distinguishes three attribution levels:
\begin{itemize}
    \item \textbf{direct attribution}, when a site visit or action can be linked to a known AI referral source or a tracked answer-click path using referral or session instrumentation;
    \item \textbf{delivery-side attribution}, when the exposed page variant and request type can be recovered from versioned page snapshots, variant identifiers, and CDN or edge delivery records even if browser analytics are incomplete;
    \item \textbf{observational correlation}, when only time-window-level associations between AI visibility and site behavior can be measured from aggregated analytics events such as visits, engagement, subscription starts, or purchases.
\end{itemize}

\section{System Architecture}

The implementation is organized as a batch-oriented recommendation and evaluation pipeline with explicit versioning of page content, prompts, answers, labels, recommendations, and aggregated metrics. It supports five main tasks:
\begin{itemize}
    \item capture and version brand-controlled pages and their revisions;
    \item compute page-level Information Gain scores and aggregate them into cluster-level static compliance signals;
    \item generate content recommendations from score breakdowns and runtime failures;
    \item execute cluster-specific prompts against selected AI systems and collect Query--Answer Traces (QAT);
    \item aggregate page-side and trace-side evidence into page-level edit targets, cluster-level metrics, recommendation priorities, and intervention decisions.
\end{itemize}

Page scoring runs on stored page snapshots, cluster scoring aggregates those page scores and cross-page signals, and runtime evaluation runs on answers collected from AI systems during scheduled measurement windows. The layers are linked by page URL, topic cluster, variant, and time window.

\begin{figure}[H]
    \centering
    \includegraphics[width=0.95\linewidth]{resources/overview.png}
    \caption{Overview of the GEO recommendation and evaluation pipeline. Brand-controlled content is first processed by the static scoring flow, which computes Information Gain scores and identifies recommendation candidates. Recommended changes are reviewed and deployed either through origin-side updates or controlled experimental variants. The deployed content is then evaluated through runtime AI-answer collection, where Query--Answer Traces are aggregated into visibility, citation, and alignment metrics. Static failures and runtime failures feed back into the recommendation loop, enabling iterative optimization across content revisions and measurement cycles.}
    \label{fig:overview-geo-pipeline}
\end{figure}

\subsection{Dashboards and Product UI Outputs}

The pipeline produces outputs for dashboards and workflow integration:
\begin{itemize}
    \item page-level IG scores, dimension breakdowns, affected-page status, and score histories;
    \item cluster-level IG scores, dimension breakdowns, and score histories;
    \item recommendation bundles with issue categories, supporting page evidence, affected pages, and suggested actions;
    \item cluster-level runtime metrics across measurement windows;
    \item failure summaries grouped by cluster, variant, and AI system;
    \item comparisons between baseline and experimental variants.
\end{itemize}

\subsection{Core Components}

\paragraph{1. Page Snapshot Service.}
This component captures each relevant page revision as a versioned artifact. For every publish event, manual import, or scheduled crawl, it stores a normalized page snapshot
\[
p = \langle u, x, m, c, v, t_{\text{snap}} \rangle,
\]
where \(u\) is the canonical URL, \(x\) is extracted content, \(m\) is metadata, \(c\) is the topic cluster, \(v\) is the exposed variant identifier, and \(t_{\text{snap}}\) is the snapshot timestamp.

\paragraph{2. IG Feature Extraction and Scoring Service.}
This component consumes versioned page snapshots and computes a page feature vector
\[
\hat{\mathbf{f}}(p,C,R_C),
\]
together with the page score
\[
\hat{g}(p \mid C,R_C).
\]
It then groups scored pages by topic cluster and computes the cluster feature vector
\[
\hat{\mathbf{F}}(C,R_C),
\]
together with the operational score
\[
\hat{G}(C \mid R_C).
\]
The service supports intent coverage, specificity, structured answerability, evidence and trust support, originality, search readiness, media and product readiness, and content risk control. Page scores are retained as first-class outputs and as evidence for each cluster score.

\paragraph{3. Recommendation Generation Service.}
This component translates page score breakdowns, cluster score breakdowns, failed card rules, and runtime failure evidence into actionable content recommendations. Its outputs may include missing-claim alerts, structure improvement suggestions, evidence-strengthening suggestions, discovery-readiness fixes, media and product-data improvements, contradiction flags, and cluster-specific rewrite guidance with affected URLs. Initial recommendation generation is rule-based and operates on the active cluster evidence set together with \(IGC_C\).

\paragraph{4. Prompt Execution and QAT Collection Service.}
This component executes prompts \(q \in Q_C\) against one or more AI systems \(s \in S\) and stores the resulting traces
\[
r = \langle t, q, s, a, D, C, v \rangle,
\]
where \(v\) records the content variant exposed during collection.

\paragraph{5. Judging and Calibration Service.}
This component assigns labels to page evidence objects, cluster evidence objects, and trace objects. For page and cluster evidence objects, it estimates rubric dimensions used in IG scoring. For trace objects, it assigns alignment judgments and failure tags from IG-F1--IG-F6.

\paragraph{6. Metric Aggregation and Card Evaluation Service.}
This component computes cluster-window metrics from QAT records and evaluates the active Information Gain Card
\[
IGC_C = \langle \Sigma_C^{\text{page}}, \Sigma_C^{\text{cluster}}, \Phi_C, \tau_C, \Gamma_C, \Omega_C \rangle.
\]

\paragraph{7. Persistent Storage Layer.}
The storage layer is logically divided into three stores:
\begin{itemize}
    \item an \textbf{artifact store} for raw page snapshots, raw answer texts, and citation extraction outputs;
    \item an \textbf{operational store} for normalized page, prompt, trace, judge, recommendation, and variant entities;
    \item a \textbf{metrics store} for aggregated IG scores, cluster-window runtime metrics, card evaluations, recommendation histories, and experiment summaries.
\end{itemize}

\subsection{Core Data Entities}

The implementation uses the following primary entities:
\begin{itemize}
    \item \texttt{PageSnapshot};
    \item \texttt{PageScore};
    \item \texttt{ClusterEvidenceSet};
    \item \texttt{ClusterScore};
    \item \texttt{RecommendationBundle};
    \item \texttt{PromptTemplate};
    \item \texttt{PromptExecution};
    \item \texttt{AnswerTrace};
    \item \texttt{JudgeResult};
    \item \texttt{ClusterMetricWindow};
    \item \texttt{IGCardConfig};
    \item \texttt{VariantAssignment}.
\end{itemize}

\subsection{Processing Flow}

The end-to-end flow is as follows:
\begin{enumerate}
    \item a page publish, revision, or scheduled crawl creates a new page snapshot;
    \item the snapshot is assigned to a topic cluster and tagged with the active variant;
    \item the IG scoring service computes page-level features and a page score;
    \item the IG scoring service groups scored pages by cluster and computes cluster-level features and scores;
    \item the recommendation service generates cluster-specific recommendation bundles and attaches affected pages as edit targets and evidence;
    \item prompts associated with the same cluster are executed during scheduled measurement windows;
    \item the resulting answers are stored as QAT records and labeled by the judging service;
    \item cluster-level metrics are aggregated for the active time window;
    \item the card evaluation service checks \(\Sigma_C^{\text{page}}\), \(\Sigma_C^{\text{cluster}}\), \(\Phi_C\), and \(\tau_C\);
    \item failing clusters or pages are routed to review, revision, controlled rollout, or re-measurement workflows.
\end{enumerate}

\subsection{Variant Routing Model}

When controlled experimentation is available, the system supports a finite set of site variants
\[
\mathcal{V}=\{v_0,v_1,\dots,v_m\},
\]
where \(v_0\) denotes the baseline and \(v_k\), \(k>0\), denote experimental content variants.

Variant routing is treated as an explicit infrastructure concern rather than as an implicit property of page content. A routing layer external to the scoring pipeline assigns a variant identifier \(v\) to the exposed page response. This layer may be implemented through CDN rules, edge middleware, alternate hostnames, or controlled internal evaluation URLs, depending on deployment constraints.

Regardless of routing mechanism, the implementation requires that:
\begin{itemize}
    \item every exposed page response be associated with a stable variant identifier;
    \item every stored page snapshot include that identifier;
    \item every QAT record include the variant against which the query was executed;
    \item aggregated metrics remain segmentable by cluster, system, window, and variant.
\end{itemize}

This routing layer supports controlled evaluation of agent-facing variants before canonical rollout. In edge-based deployments, content transformations may be applied at the CDN layer without changing the origin CMS response. Every exposed response, QAT record, recommendation outcome, and downstream site event must remain segmentable by cluster, measurement window, and variant.

\subsection{Routing Constraints and Safeguards}

Variant routing is only considered valid for evaluation if the following constraints hold:
\begin{itemize}
    \item the assigned variant is logged consistently in both page-side and QAT-side records;
    \item prompt execution targets a known variant deterministically;
    \item routing does not create ambiguous attribution between baseline and experimental content;
    \item inactive or deprecated variants cannot be selected by scheduled QAT jobs;
    \item the evaluation system can reconstruct which exact content version was exposed during each measurement window.
\end{itemize}

\subsection{Operational Scope for V1}

The first implementation cycle makes the following simplifying decisions:
\begin{itemize}
    \item page scoring and cluster aggregation are performed offline on versioned page snapshots;
    \item runtime evaluation is performed on scheduled prompt batches rather than live user traffic;
    \item the prompt inventory is human-authored and versioned;
    \item before--after comparison is the primary evaluation design;
    \item adversarial prompt exploration is out of scope for the initial deployment;
    \item variant routing, if used, is limited to controlled evaluation scenarios with deterministic tagging.
\end{itemize}

\section{Implementation Decisions}

\paragraph{V1.}

The initial implementation uses the broad cluster \texttt{unknown} for all captured Frescopa pages and content. The static scorer is two-level: page snapshots receive local scores, and scored pages are aggregated into cluster-level scores for topic evaluation, runtime comparison, and reporting. The dimension set includes intent coverage, specificity, structured answerability, trust support, originality, search readiness, media readiness, and risk control.

Scoring is heuristic and deterministic, with human or LLM-based judging reserved for later calibration. Runtime evaluation uses presence, prominence, citation share, and alignment. The operational output is a recommendation workflow that maps low-performing page and cluster dimensions through \(\Omega_C\) to cluster-specific recommendations with affected pages.

\paragraph{Implementation references.}

The following search-system references inform implementation details while the model remains search-engine-agnostic:
\begin{itemize}
    \item \url{https://developers.google.com/search/docs/fundamentals/seo-starter-guide}
    \item \url{https://developers.google.com/search/docs/fundamentals/how-search-works}
    \item \url{https://developers.google.com/search/docs/fundamentals/creating-helpful-content}
    \item \url{https://developers.google.com/search/docs/fundamentals/ai-optimization-guide}
    \item \url{https://developers.google.com/search/docs/fundamentals/using-gen-ai-content}
\end{itemize}

\bibliographystyle{alpha}
% \bibliography{sample}

\end{document}
